\section{Data}
\label{sec:data}

The series come from one static administrative extract, labeled \texttt{extracao=2026-03}.
The extract is a research snapshot of enrolled-debt stock and of collections recorded on state tax payment slips (GARE).
It is not a live stream.
If a source file is absent, that grain is skipped.
We do not impute a series that was not in the extract.

\paragraph{Enrolled debt and collection.}
Enrolled debt is the stock of tax debts formally registered for collection after the administrative assessment.
The monthly panel (\texttt{painel\_mensal\_estoque\_arrecadacao}) pairs that stock with realized collection.
In this freeze the panel has $n = 119$ months and 14 columns.
The monthly target is total collection, \texttt{valor\_total}.
Tabular models on this grain use lags 1, 3, and 12 of that collection series and the same lags of the stock field \texttt{valor\_sem\_honorarios}.
The short length is material: a 12-month test block leaves limited train mass for any high-capacity fit.

\paragraph{Calendar-daily payment slips (Ranking~A).}
The daily source is \texttt{gare\_diario\_2026-03}, an aggregate of payment slips.
Each civil day carries the sum of slip values, \texttt{valor\_sum}, and the slip count, \texttt{n\_guias}.
We reindex the file to a continuous daily calendar and set missing days to zero on both fields.
The forecast target is \texttt{valor\_sum}.
Zeros are part of the target.
They mark weekends and days with no slip, which are ordinary in this collection system.
Ranking~A scores every calendar day, zeros included.
A side score on days with positive collection, sMAPE$_{y \gt 0}$, describes active days only.
It does not replace Ranking~A.

\paragraph{Business-day ablation (Ranking~B).}
Ranking~B restricts the same slip aggregate to Monday--Friday.
In this freeze that series has $n = 2672$ days.
Brazilian national holidays remain in the series.
The ablation aligns the time index with a business-day frequency and does not apply a holiday regressor.
Ranking~B is a different target, with its own train, validation, and test blocks.
It is not a filter applied after the fact to the calendar-daily holdout.

\paragraph{Weekly aggregate.}
We sum daily \texttt{valor\_sum} (and the slip count) by week ending Sunday.
The seasonal reference lag is 52, one year of weeks.
Sunday close keeps the week aligned with the civil calendar used in Ranking~A.
We do not use an ISO week starting Monday.

\paragraph{Covariates by grain.}
All four grains use the same extract and the same metric code.
They differ in sampling frequency, in the seasonal lag of the reference forecast, and in the tabular covariate set.
Monthly tabular models use the lags named above.
Calendar-daily tabular models use lags 1, 7, 14, and 28 of \texttt{valor\_sum} and of \texttt{n\_guias}, plus day-of-week and month.
Business-day tabular models use lags 1, 5, 10, and 20 in business-day steps of the same two series, plus day-of-week and month.
Weekly tabular models use lags 1, 4, 13, and 52 of the weekly sum and of the weekly slip count, plus week-of-year and month.
The seasonal naive, SARIMAX, Prophet, the perceptron, TimesFM, and the generative samplers see the target history (Prophet also sees the time index that carries its seasonal terms).
Gaussian processes and kernel ridge use the tabular covariates of the grain, not the target alone.
